%% 1-21 広さを くらべよう：面積（10ページ、p.212〜221）　第1巻の最後の章
%% 第1巻のきまり：15 pt・行送り28 pt、分かち書き（文節の間に半角の空白）、漢字はすべてルビ、式は全角の＋＝で書く
%% 小4の内容を、小3向けに早めにやさしく扱う。辺・正方形・長方形は 1-19、かけ算は 1-12 と同じ ことばにそろえる。

\makeatletter
% 形や部屋につける名前の札（1-19 と同じ）  \bkqttag{(x,y)}{A}
\newcommand\bkqttag[2]{\node[circle, draw=bkInk!70, fill=white, line width=0.7pt, text=bkInk,
  font=\sffamily\bfseries\footnotesize, minimum size=5.2mm, inner sep=0pt] at #1 {#2};}
% マスでできた形  \bkqtcells{(x,y)}{1マス}{色}{よこ番号/たて番号,…}
\newcommand\bkqtcells[4]{\foreach \bkqt@i/\bkqt@j in {#4} {\bkcell{#1}{\bkqt@i}{\bkqt@j}{#2}{#3}}}
% へやのゆか（上から見た図。木のゆか）  \bkqtfloor{(左下)}{はば}{たかさ}
\newcommand\bkqtfloor[3]{\begin{scope}[shift={#1}]
  \filldraw[fill=bkYellow!28, draw=bkStem, line width=1.2pt] (0,0) rectangle (#2,#3);
  \foreach \bkqt@y in {2.5,5,...,#3} {\ifdim\bkqt@y pt<#3pt\relax\draw[bkStem!35, line width=0.4pt] (0.6,\bkqt@y) -- (#2-0.6,\bkqt@y);\fi}
\end{scope}}
% たたみ1まい  \bkqttatami{(左下)}{はば}{たかさ}
\newcommand\bkqttatami[3]{\filldraw[fill=bkLeaf!22, draw=bkStem!80, line width=0.8pt] #1 rectangle ++(#2,#3);}
\makeatother

\bkchapter{1}{21}{\R{広}{ひろ}さを くらべよう：\R{面積}{めん|せき}}{\bklevel{1}}{ひろさを くらべよう：めんせき}
\bkmichishirube{1-8,1-12,1-19}{2-3,4-5,5-10}

\begin{goal}
\begin{itemize}[label={\color{bkTheme}●}, leftmargin=1.4\zw, itemsep=0pt, topsep=0pt]
\item \R{広}{ひろ}さを、 \R{同}{おな}じ \R{大}{おお}きさの マスの \R{数}{かず}で くらべられる。
\item 1\,cm\textsuperscript{2} を もとに、 \R{面積}{めん|せき}を \R{数}{かず}で あらわせる。
\item \R{長方形}{ちょう|ほう|けい}の \R{面積}{めん|せき}を たて × よこ で もとめられる。
\end{itemize}
\end{goal}

\begin{hatena}
\begin{bkfigure}
\begin{tikzpicture}[x=1mm, y=1mm]
  % A の へや
  \fill[white, rounded corners=3mm] (0,0) rectangle (39,54);
  \node[circle, fill=bkYellow, text=white, font=\sffamily\bfseries, minimum size=6.5mm, inner sep=0pt] at (6,48) {A};
  \node[anchor=west, font=\small, text=bkInk] at (10.5,48) {A の へや};
  \bkqtfloor{(3.5,11)}{32}{19.2}
  % B の へや
  \begin{scope}[shift={(41.5,0)}]
  \fill[white, rounded corners=3mm] (0,0) rectangle (39,54);
  \node[circle, fill=bkYellow, text=white, font=\sffamily\bfseries, minimum size=6.5mm, inner sep=0pt] at (6,48) {B};
  \node[anchor=west, font=\small, text=bkInk] at (10.5,48) {B の へや};
  \bkqtfloor{(6.7,8)}{25.6}{25.6}
  \end{scope}
  % かさねると
  \begin{scope}[shift={(83,0)}]
  \fill[white, rounded corners=3mm] (0,0) rectangle (39,54);
  \node[anchor=west, font=\small, text=bkInk] at (3,48) {かさねると？};
  \fill[bkRoute4!30] (3.5,9) rectangle (29.1,34.6);
  \fill[bkOrange!35, opacity=0.75] (3.5,9) rectangle (35.5,28.2);
  \draw[bkRoute4!75!black, line width=1.2pt] (3.5,9) rectangle (29.1,34.6);
  \draw[bkOrange!90!black, line width=1.2pt, dash pattern=on 3pt off 1.5pt] (3.5,9) rectangle (35.5,28.2);
  \node[font=\sffamily\bfseries\footnotesize, text=bkRoute4!60!black] at (16.3,31.4) {B};
  \node[font=\sffamily\bfseries\footnotesize, text=bkOrange!80!black] at (32.3,18.6) {A};
  \node[font=\sffamily\bfseries\footnotesize, text=bkInk] at (19.5,4.5) {どちらも はみ\R{出}{だ}す};
  \end{scope}
\end{tikzpicture}
\end{bkfigure}
\noindent A と B、 どっちの へやが \R{広}{ひろ}いかな？\\
かさねて みても、 どちらも はみ\R{出}{だ}して しまう。\\
どう したら くらべられるかな？
\end{hatena}

\clearpage
%% ---------------- p.213 やってみよう ----------------
\begin{yattemiyou}[かさねて くらべよう]
\begin{enumerate}[label={\bkcircled[bkGreen]{\arabic*}}, leftmargin=2.2\zw, itemsep=2pt, topsep=0pt]
\item ハンカチと カードを、 かどを そろえて かさねよう。 はみ\R{出}{だ}した ほうが \R{広}{ひろ}い。
\end{enumerate}
\begin{bkfigure}
\begin{tikzpicture}[x=1mm, y=1mm]
  \fill[white, rounded corners=3mm] (0,0) rectangle (122,30);
  \filldraw[fill=bkRoute5!25, draw=bkRoute5!70!black, line width=1pt] (34,4) rectangle (56,26);
  \foreach \k in {1,2,3} {\draw[bkRoute5!50, line width=0.6pt] ({34+5.5*\k},4) -- ({34+5.5*\k},26) (34,{4+5.5*\k}) -- (56,{4+5.5*\k});}
  \filldraw[fill=white, draw=bkPurple, line width=1pt] (34,4) rectangle (48,15);
  \fill[bkPurple!60] (37,11) rectangle (45,12.5);\fill[bkPurple!30] (37,7.5) rectangle (43,9);
  \node[anchor=west, font=\sffamily\bfseries\small, text=bkGreen!70!black, align=left] at (62,15) {ハンカチが はみ\R{出}{だ}した。\\ハンカチの ほうが \R{広}{ひろ}い！};
\end{tikzpicture}
\end{bkfigure}
\begin{enumerate}[label={\bkcircled[bkGreen]{\arabic*}}, leftmargin=2.2\zw, itemsep=2pt, topsep=4pt, start=2]
\item きいろと みどりの \R{紙}{かみ}は、 かさねると どちらも はみ\R{出}{だ}します。 \R{同}{おな}じ \R{大}{おお}きさの シールを はって、 \R{何}{なん}まい はれるか かぞえよう。
\end{enumerate}
\begin{bkfigure}
\begin{tikzpicture}[x=1mm, y=1mm]
  \fill[white, rounded corners=3mm] (0,0) rectangle (122,36);
  % きいろ（よこ 4まい × たて 2まい）。 マスの 線は 紙の かどから 7\,mm ずつ
  \begin{scope}[shift={(18,8)}]
  \fill[bkYellow!35] (0,0) rectangle (28,14);
  \draw[bkYellow!70!black, line width=0.5pt] (0,0) grid[step=7] (28,14);
  \bkcell{(0,0)}{3}{0}{7}{bkOrange!70}
  \draw[bkYellow!70!black, line width=1.2pt] (0,0) rectangle (28,14);
  \end{scope}
  \node[anchor=base, font=\sffamily\bfseries\footnotesize] at (32,28.5) {きいろ};
  % みどり（3まい × 3まい）
  \begin{scope}[shift={(83,4)}]
  \fill[bkRoute4!25] (0,0) rectangle (21,21);
  \draw[bkRoute4!70!black, line width=0.5pt] (0,0) grid[step=7] (21,21);
  \draw[bkRoute4!70!black, line width=1.2pt] (0,0) rectangle (21,21);
  \end{scope}
  \node[anchor=base, font=\sffamily\bfseries\footnotesize] at (93.5,28.5) {みどり};
  \node[anchor=west, font=\sffamily\footnotesize, text=bkGray, align=left] at (50.5,13) {シールを\\1まい\\はった ところ};
  \draw[bkGray, line width=0.6pt, -{Stealth[length=3pt]}] (50.3,11.5) -- (44.2,11.5);
\end{tikzpicture}
\end{bkfigure}
\noindent\hspace*{1\zw}きいろ\enspace\bkkakikomi\enspace まい\hspace{2\zw}みどり\enspace\bkkakikomi\enspace まい\\
\hspace*{1\zw}\R{広}{ひろ}いのは\enspace\bkkakikomi[24mm]\par
\noindent\hspace*{1\zw}{\small\color{bkOrange!58!black}ヒント：かぞえた マスに しるしを つけると、 まちがえにくいよ。}
\end{yattemiyou}

\bkfukidashi{mina}{きいろは 8まい、 みどりは 9まい。 みどりが \R{広}{ひろ}いね！\\マスの \R{数}{かず}を くらべれば いいんだ。}

\clearpage
%% ---------------- p.214〜215 数学の道具 ----------------
\bksection{doguu}{\R{数学}{すう|がく}の \R{道具}{どう|ぐ}}

\bksubsection{1}{マスの \R{数}{かず}で くらべる}

\label{bk:1-21:masu}\bkhatenaref の へやに、 \R{同}{おな}じ \R{大}{おお}きさの マスを ならべて みよう。

\begin{bkfigure}
\begin{tikzpicture}[x=1mm, y=1mm]
  \begin{scope}[shift={(10,0)}]\fill[bkYellow!28] (0,0) rectangle (35,21);\draw[bkStem!70, line width=0.6pt] (0,0) grid[step=7] (35,21);\draw[bkStem, line width=1.2pt] (0,0) rectangle (35,21);\end{scope}
  \bkqttag{(5,24)}{A}
  \node[anchor=north, font=\sffamily\bfseries\small] at (27.5,-1.5) {A　15マス};
  \begin{scope}[shift={(78,0)}]\fill[bkYellow!28] (0,0) rectangle (28,28);\draw[bkStem!70, line width=0.6pt] (0,0) grid[step=7] (28,28);\draw[bkStem, line width=1.2pt] (0,0) rectangle (28,28);\end{scope}
  \bkqttag{(73,31)}{B}
  \node[anchor=north, font=\sffamily\bfseries\small] at (92,-1.5) {B　16マス};
\end{tikzpicture}
\end{bkfigure}

\noindent A は たて 3マス、 よこ 5マス。 B は たて 4マス、 よこ 4マス。 1つずつ かぞえると、 A は 15マス、 B は 16マス です。\\
\bkhatenaref の \R{答}{こた}え：16 － 15 ＝ 1で、 B の へやが 1マス \R{広}{ひろ}い。

\par\vspace{4pt}
\noindent\begin{minipage}[c]{0.56\linewidth}
ただし、 マスの \R{大}{おお}きさが ちがうと、 くらべられません。
\R{右}{みぎ}の C は \R{大}{おお}きい マスで 4こ、 D は \R{小}{ちい}さい マスで 9こ。
マスの \R{数}{かず}は D が \R{多}{おお}いのに、 \R{広}{ひろ}いのは C です。
だから、 かならず \R{同}{おな}じ \R{大}{おお}きさの マスで かぞえます。
\end{minipage}\hfill
\begin{minipage}[c]{0.42\linewidth}\centering
\begin{tikzpicture}[x=1mm, y=1mm]
  \fill[bkYellow!35] (0,0) rectangle (20,20);
  \draw[bkStem!80, line width=0.6pt] (0,0) grid[step=10] (20,20);
  \draw[bkStem, line width=1.1pt] (0,0) rectangle (20,20);
  \bkqttag{(10,24.5)}{C}
  \node[anchor=north, font=\sffamily\bfseries\footnotesize] at (10,-1.2) {4マス};
  \begin{scope}[shift={(28,2)}]
  \fill[bkRoute4!30] (0,0) rectangle (15,15);
  \draw[bkRoute4!70!black, line width=0.6pt] (0,0) grid[step=5] (15,15);
  \draw[bkRoute4!70!black, line width=1.1pt] (0,0) rectangle (15,15);
  \bkqttag{(7.5,22.5)}{D}
  \node[anchor=north, font=\sffamily\bfseries\footnotesize] at (7.5,-3.2) {9マス};
  \end{scope}
\end{tikzpicture}
\end{minipage}

\bksubsection{2}{1\,cm\textsuperscript{2}（1へいほうセンチメートル）}

\label{bk:1-21:cm}\R{長}{なが}さを はかる ときは、 だれでも \R{同}{おな}じ 1\,cm を つかいました（1-8）。
\R{広}{ひろ}さも、 みんなで \R{同}{おな}じ マスを つかえば、 \R{数}{かず}で つたえられます。
そこで、 1\R{辺}{へん}が 1\,cm の \R{正方形}{せい|ほう|けい}の \R{広}{ひろ}さを もとに して、 1\,cm\textsuperscript{2} と かきます。
\R{右上}{みぎ|うえ}の \R{小}{ちい}さな 2は、 たてと よこの 2つの むきに ひろがった \R{広}{ひろ}さの しるしです。
\R{広}{ひろ}さの ことを \R{面積}{めん|せき}と いいます。\bkidx{めんせき}{\R{面積}{めん|せき}}\bkidx{へいほうせんちめーとる}{\R{平方}{へい|ほう}センチメートル（cm\textsuperscript{2}）}

\bkkigocards{\bkkigocard{cm\textsuperscript{2}}{へいほう\\センチメートル}{\R{面積}{めん|せき}の たんい}{square\\centimeter}{1-21}}{%
\begin{tikzpicture}[x=1mm, y=1mm]
  \bkqtcells{(0,31)}{10}{bkYellow!60}{0/0,1/0,2/0,3/0}
  \node[anchor=west, font=\sffamily\bfseries\footnotesize] at (42,36) {4\,cm\textsuperscript{2}};
  \bkqtcells{(0,8)}{10}{bkRoute4!40}{0/0,1/0,0/1,1/1}
  \node[anchor=west, font=\sffamily\bfseries\footnotesize] at (22,23) {4\,cm\textsuperscript{2}};
  \node[anchor=west, font=\sffamily\footnotesize, align=left] at (22,13) {1マスが 1\,cm\textsuperscript{2}\\（ほんとうの \R{大}{おお}きさ）};
  \node[anchor=west, font=\sffamily\footnotesize] at (0,2.5) {\R{形}{かたち}が ちがっても \R{面積}{めん|せき}は \R{同}{おな}じ};
\end{tikzpicture}}

\clearpage
\bksubsection{3}{\R{長方形}{ちょう|ほう|けい}の \R{面積}{めん|せき}は かけ\R{算}{ざん}で}

\label{bk:1-21:kakezan}マスを 1つずつ かぞえなくても、 かけ\R{算}{ざん}で もとめられます。\par\vspace{3pt}

\noindent\begin{minipage}[c]{0.42\linewidth}\centering
\begin{tikzpicture}[x=1mm, y=1mm]
  \bkgrid{(0,0)}{4}{3}{8}
  \draw[bkTheme, line width=1.2pt] (0,0) rectangle (32,24);
  \foreach \i in {0,...,3} {\foreach \j in {0,1,2} {\bkohajiki{(4+8*\i,4+8*\j)}{2.4}{bkRoute2}}}
  \node[anchor=east, font=\sffamily\bfseries\footnotesize] at (-1,12) {たて 3\,cm};
  \node[anchor=north, font=\sffamily\bfseries\footnotesize] at (16,-1) {よこ 4\,cm};
\end{tikzpicture}
\end{minipage}\hfill
\begin{minipage}[c]{0.56\linewidth}
1マスの 1\R{辺}{へん}は 1\,cm。\\
だから たて 3\,cm には マスが 3こ、 よこ 4\,cm には 4れつ ならびます。\\
1れつ 3こが 4れつで、 1-12の 「1つ\R{分}{ぶん} × いくつ\R{分}{ぶん}」。\\
3 × 4 ＝ 12 で、 12\,cm\textsuperscript{2} です。
\end{minipage}

\noindent 1つずつ かぞえても 12こ。 かけ\R{算}{ざん}と \R{同}{おな}じ \R{答}{こた}えに なります。
よこ × たて の 4 × 3 でも 12（1-13の 「\R{入}{い}れかえても \R{同}{おな}じ」）です。

\begin{taisetsu}
\R{長方形}{ちょう|ほう|けい}の \R{面積}{めん|せき} ＝ たて × よこ\\
\R{正方形}{せい|ほう|けい}も、 たて × よこ で もとめられる。
\end{taisetsu}

\bksubsection{4}{\R{大}{おお}きな \R{広}{ひろ}さ：1\,m\textsuperscript{2}}

\label{bk:1-21:m}\R{教室}{きょう|しつ}の ゆかを 1\,cm\textsuperscript{2} の マスで かぞえると、 マスが \R{多}{おお}すぎて たいへんです。
\R{長}{なが}さで、 \R{長}{なが}い ものは cm ではなく m で はかったのと \R{同}{おな}じ（1-8）です。
\R{大}{おお}きな \R{広}{ひろ}さは、 \R{大}{おお}きな マス 1\,m\textsuperscript{2} で かぞえます。

\noindent 1\R{辺}{へん}が 1\,m の \R{正方形}{せい|ほう|けい}の \R{広}{ひろ}さを、 1\,m\textsuperscript{2} と かきます。\bkidx{へいほうめーとる}{\R{平方}{へい|ほう}メートル（m\textsuperscript{2}）}

\bkkigocards{\bkkigocard{m\textsuperscript{2}}{へいほう メートル}{\R{面積}{めん|せき}の たんい}{square meter}{1-21}}{%
\begin{tikzpicture}[x=1mm, y=1mm]
  \filldraw[fill=bkGray!12, draw=bkInk!60, line width=0.8pt] (0,0) -- (40,0) -- (52,16) -- (12,16) -- cycle;
  \foreach \t in {0.25,0.5,0.75} {\draw[bkGray!50, line width=0.5pt] ({40*\t},0) -- ({12+40*\t},16);}
  \draw[bkGray!50, line width=0.5pt] (6,8) -- (46,8);
  \bkkid{(19,4)}{18}{bkRoute6}{bkHair}\bkkid{(33,5)}{17}{bkRoute4}{bkHair}
  \node[anchor=north, font=\sffamily\bfseries\footnotesize] at (20,-0.5) {1\,m};
  \node[anchor=west, font=\sffamily\bfseries\footnotesize] at (47,7) {1\,m};
  \node[anchor=north, font=\sffamily\footnotesize, text=bkGray, align=left] at (26,-5) {\R{新聞紙}{しん|ぶん|し}（ひらいた もの）\\2まいと すこしの \R{広}{ひろ}さ};
\end{tikzpicture}}

\par\vspace{3pt}
\noindent \R{教室}{きょう|しつ}は たて 7\,m、 よこ 9\,m ほど。 1\,m\textsuperscript{2} の マスが、 たてに 7こ ならぶ れつが 9れつ。
7 × 9 ＝ 63 で、 だいたい 63\,m\textsuperscript{2} です。

\begin{column}{1\,m\textsuperscript{2} は 1\,cm\textsuperscript{2} の いくつ\R{分}{ぶん}？}
1\,m は 100\,cm（1-8）。 1\,m\textsuperscript{2} の \R{正方形}{せい|ほう|けい}には、 1\,cm\textsuperscript{2} の マスが たてに 100こ、 よこに 100れつ ならびます。
100が 100こで 10000こ（1-2）。 1\,m\textsuperscript{2} は 1\,cm\textsuperscript{2} の 10000こ\R{分}{ぶん}です。 \R{教室}{きょう|しつ}を cm\textsuperscript{2} で かぞえると たいへんな わけが わかるね。
\end{column}

\par\vspace{6pt}
%% ---------------- p.216〜217 例題 ----------------
\begin{reidai}[マスで くらべる]
A と B、 どちらが \R{何}{なん}マス \R{広}{ひろ}い？\par
\noindent\begin{minipage}[c]{0.52\linewidth}
\begin{tikzpicture}[x=1mm, y=1mm]
  \bkgrid{(0,0)}{11}{3}{5}
  \bkqtcells{(0,0)}{5}{bkYellow!60}{0/0,1/0,2/0,0/1,1/1,2/1,0/2}
  \bkqtcells{(0,0)}{5}{bkRoute4!40}{6/0,7/0,8/0,9/0,6/1,7/1,8/1,9/1}
  \bkqttag{(7.5,12.5)}{A}\bkqttag{(40,12.5)}{B}
\end{tikzpicture}
\end{minipage}\hfill
\begin{minipage}[c]{0.46\linewidth}
\kangaekata A 7マス、\\
B 8マス。
\end{minipage}\par
\shiki 8 － 7 ＝ 1 \kotaenext B が 1マス \R{広}{ひろ}い
\end{reidai}

\begin{jibun}
\noindent\begin{minipage}[c]{0.52\linewidth}
\begin{tikzpicture}[x=1mm, y=1mm]
  \bkgrid{(0,0)}{10}{3}{5}
  \bkqtcells{(0,0)}{5}{bkYellow!60}{0/0,1/0,2/0,3/0,4/0,0/1,1/1,2/1,3/1,4/1}
  \bkqtcells{(0,0)}{5}{bkRoute4!40}{6/0,7/0,8/0,6/1,7/1,8/1,6/2,7/2}
  \bkqttag{(12.5,12.5)}{A}\bkqttag{(45,12.5)}{B}
\end{tikzpicture}
\end{minipage}\hfill
\begin{minipage}[c]{0.46\linewidth}
A と B、 どちらが\\\R{何}{なん}マス \R{広}{ひろ}い？
\end{minipage}
\bkshikiran[10\zw]
\bkjans{10 － 8 ＝ 2\quad A が 2マス \R{広}{ひろ}い}
\end{jibun}

\begin{reidai}[cm\textsuperscript{2} で あらわす]
この \R{形}{かたち}の \R{面積}{めん|せき}は \R{何}{なん}\,cm\textsuperscript{2}？\par
\noindent\begin{minipage}[c]{0.52\linewidth}
\begin{tikzpicture}[x=1mm, y=1mm]
  \bkgrid{(0,0)}{8}{3}{5}
  \bkqtcells{(0,0)}{5}{bkRoute6!40}{1/2,2/2,3/2,2/1,2/0,1/0,3/0}
  \node[anchor=west, font=\sffamily\footnotesize, text=bkGray, align=left] at (42,7.5) {1マスは\\1\,cm\textsuperscript{2}};
\end{tikzpicture}
\end{minipage}\hfill
\begin{minipage}[c]{0.46\linewidth}
\kangaekata\\1\,cm\textsuperscript{2} の マスが 7こ。
\end{minipage}\par
\kotae 7\,cm\textsuperscript{2}
\end{reidai}

\begin{jibun}
この \R{形}{かたち}の \R{面積}{めん|せき}は \R{何}{なん}\,cm\textsuperscript{2}？\par
\noindent\begin{minipage}[c]{0.52\linewidth}
\begin{tikzpicture}[x=1mm, y=1mm]
  \bkgrid{(0,0)}{8}{3}{5}
  \bkqtcells{(0,0)}{5}{bkRoute5!35}{0/0,1/0,2/0,3/0,0/1,1/1,2/1,0/2,1/2}
  \node[anchor=west, font=\sffamily\footnotesize, text=bkGray, align=left] at (42,7.5) {1マスは\\1\,cm\textsuperscript{2}};
\end{tikzpicture}
\end{minipage}\hfill
\begin{minipage}[c]{0.46\linewidth}
\bkkotaeran[24mm]
\end{minipage}
\bkjans{9\,cm\textsuperscript{2}}
\end{jibun}

\begin{reidai}[\R{長方形}{ちょう|ほう|けい}の \R{面積}{めん|せき}]
たて 4\,cm、 よこ 6\,cm の \R{長方形}{ちょう|ほう|けい}の \R{面積}{めん|せき}は？\par
\noindent\begin{minipage}[c]{0.56\linewidth}
\kangaekata 1\,cm\textsuperscript{2} の マスが、 たてに 4こ ならぶ れつが 6れつ。\\
\R{長方形}{ちょう|ほう|けい}の \R{面積}{めん|せき} ＝ たて × よこ
\end{minipage}\hfill
\begin{minipage}[c]{0.42\linewidth}\centering
\begin{tikzpicture}[x=1mm, y=1mm]
  \fill[bkRoute6!20] (0,0) rectangle (30,20);
  \draw[bkRoute6!60, line width=0.5pt] (0,0) grid[step=5] (30,20);
  \draw[bkRoute6!80!black, line width=1.2pt] (0,0) rectangle (30,20);
  \node[anchor=east, font=\sffamily\bfseries\footnotesize] at (-1,10) {4\,cm};
  \node[anchor=north, font=\sffamily\bfseries\footnotesize] at (15,-1) {6\,cm};
\end{tikzpicture}
\end{minipage}\par
\shiki 4 × 6 ＝ 24 \kotaenext 24\,cm\textsuperscript{2}
\end{reidai}

\begin{jibun}
たて 3\,m、 よこ 5\,m の へやの \R{面積}{めん|せき}は \R{何}{なん}\,m\textsuperscript{2}？
\bkshikiran
\bkjans{3 × 5 ＝ 15\quad 15\,m\textsuperscript{2}}
\end{jibun}

\begin{reidai}[\R{正方形}{せい|ほう|けい}の \R{面積}{めん|せき}]
1\R{辺}{へん}が 5\,cm の \R{正方形}{せい|ほう|けい}の \R{面積}{めん|せき}は？\par
\noindent\begin{minipage}[c]{0.58\linewidth}
\kangaekata \R{正方形}{せい|ほう|けい}は、 たても よこも 5\,cm（1-19）。
\R{長方形}{ちょう|ほう|けい}と \R{同}{おな}じように、 たて × よこ で もとめる。
\end{minipage}\hfill
\begin{minipage}[c]{0.4\linewidth}\centering
\begin{tikzpicture}[x=1mm, y=1mm]
  \fill[bkRoute2!20] (0,0) rectangle (20,20);
  \draw[bkRoute2!60, line width=0.5pt] (0,0) grid[step=4] (20,20);
  \draw[bkRoute2!80!black, line width=1.2pt] (0,0) rectangle (20,20);
  \node[anchor=east, font=\sffamily\bfseries\footnotesize] at (-1,10) {5\,cm};
  \node[anchor=north, font=\sffamily\bfseries\footnotesize] at (10,-1) {5\,cm};
\end{tikzpicture}
\end{minipage}\par
\shiki 5 × 5 ＝ 25 \kotaenext 25\,cm\textsuperscript{2}
\end{reidai}

\begin{jibun}
1\R{辺}{へん}が 6\,cm の \R{正方形}{せい|ほう|けい}の \R{面積}{めん|せき}は \R{何}{なん}\,cm\textsuperscript{2}？
\bkshikiran
\bkjans{6 × 6 ＝ 36\quad 36\,cm\textsuperscript{2}}
\end{jibun}

\begin{reidai}[\R{面積}{めん|せき}から \R{辺}{へん}の \R{長}{なが}さを もとめる]
\R{面積}{めん|せき}が 12\,cm\textsuperscript{2}、 たてが 3\,cm の \R{長方形}{ちょう|ほう|けい}。 よこは \R{何}{なん}cm？\par
\noindent\begin{minipage}[c]{0.6\linewidth}
\kangaekata たて × よこ ＝ \R{面積}{めん|せき} なので、 3 × □ ＝ 12。\\
3の だんで 12に なるのは 3 × 4（1-13）。\\
わり\R{算}{ざん}なら 12 ÷ 3 ＝ 4（1-14）。
\end{minipage}\hfill
\begin{minipage}[c]{0.38\linewidth}\centering
\begin{tikzpicture}[x=1mm, y=1mm]
  \fill[bkRoute5!25] (0,0) rectangle (20,15);
  \draw[bkRoute5!60, line width=0.5pt] (0,0) grid[step=5] (20,15);
  \draw[bkRoute5!80!black, line width=1.2pt] (0,0) rectangle (20,15);
  \node[anchor=east, font=\sffamily\bfseries\footnotesize] at (-1,7.5) {3\,cm};
  \node[anchor=north, font=\sffamily\bfseries\footnotesize] at (10,-1) {？\,cm};
\end{tikzpicture}
\end{minipage}\par
\kotae 4\,cm\quad（たしかめ：3 × 4 ＝ 12）
\end{reidai}

\begin{jibun}
\R{面積}{めん|せき}が 20\,cm\textsuperscript{2}、 たてが 4\,cm の \R{長方形}{ちょう|ほう|けい}。 よこは \R{何}{なん}cm？
\bkshikiran
\bkjans{4 × □ ＝ 20、 4 × 5 ＝ 20（20 ÷ 4 ＝ 5）\quad 5\,cm}
\end{jibun}

\begin{reidai}[\R{分}{わ}けて たす]
この \R{形}{かたち}の \R{面積}{めん|せき}は \R{何}{なん}\,cm\textsuperscript{2}？\par
\noindent\begin{minipage}[c]{0.4\linewidth}\centering
\begin{tikzpicture}[x=1mm, y=1mm]
  \fill[bkRoute4!35] (0,0) -- (24,0) -- (24,8) -- (12,8) -- (12,16) -- (0,16) -- cycle;
  \bkgrid{(0,0)}{6}{4}{4}
  \draw[bkRoute4!70!black, line width=1.2pt, line join=round] (0,0) -- (24,0) -- (24,8) -- (12,8) -- (12,16) -- (0,16) -- cycle;
  \draw[bkTheme, line width=1.1pt, dash pattern=on 2.5pt off 1.5pt] (0,8) -- (12,8);
  \node[anchor=south, font=\sffamily\bfseries\footnotesize, inner sep=1pt] at (6,16.3) {3\,cm};
  \node[anchor=north, font=\sffamily\bfseries\footnotesize, inner sep=1pt] at (12,-0.3) {6\,cm};
  \node[anchor=east, font=\sffamily\bfseries\footnotesize, inner sep=1pt] at (-0.5,12) {2\,cm};
  \node[anchor=east, font=\sffamily\bfseries\footnotesize, inner sep=1pt] at (-0.5,4) {2\,cm};
\end{tikzpicture}
\end{minipage}\hfill
\begin{minipage}[c]{0.58\linewidth}
\kangaekata 2つに \R{分}{わ}ける。\\
\R{上}{うえ}　2 × 3 ＝ 6\\
\R{下}{した}　2 × 6 ＝ 12
\end{minipage}\par
\shiki 6 ＋ 12 ＝ 18 \kotaenext 18\,cm\textsuperscript{2}

\noindent\begin{minipage}[c]{0.4\linewidth}\centering
\begin{tikzpicture}[x=1mm, y=1mm]
  \fill[bkRoute4!35] (0,0) -- (24,0) -- (24,8) -- (12,8) -- (12,16) -- (0,16) -- cycle;
  \bkgrid{(0,0)}{6}{4}{4}
  \draw[bkRoute4!70!black, line width=1.2pt, line join=round] (0,0) -- (24,0) -- (24,8) -- (12,8) -- (12,16) -- (0,16) -- cycle;
  \draw[bkTheme, line width=1.1pt, dash pattern=on 2.5pt off 1.5pt] (12,0) -- (12,8);
\end{tikzpicture}
\end{minipage}\hfill
\begin{minipage}[c]{0.58\linewidth}
たてに \R{分}{わ}けても できます。\\
\R{左}{ひだり}　4 × 3 ＝ 12\\
\R{右}{みぎ}　2 × 3 ＝ 6\\
12 ＋ 6 ＝ 18。 \R{同}{おな}じ 18\,cm\textsuperscript{2} に なります。
\end{minipage}
\end{reidai}

\begin{jibun}
\noindent\begin{minipage}[c]{0.4\linewidth}\centering
\begin{tikzpicture}[x=1mm, y=1mm]
  \fill[bkRoute5!30] (0,0) -- (20,0) -- (20,8) -- (8,8) -- (8,16) -- (0,16) -- cycle;
  \bkgrid{(0,0)}{5}{4}{4}
  \draw[bkRoute5!70!black, line width=1.2pt, line join=round] (0,0) -- (20,0) -- (20,8) -- (8,8) -- (8,16) -- (0,16) -- cycle;
  \node[anchor=south, font=\sffamily\bfseries\footnotesize, inner sep=1pt] at (4,16.3) {2\,cm};
  \node[anchor=north, font=\sffamily\bfseries\footnotesize, inner sep=1pt] at (10,-0.3) {5\,cm};
  \node[anchor=east, font=\sffamily\bfseries\footnotesize, inner sep=1pt] at (-0.5,12) {2\,cm};
  \node[anchor=east, font=\sffamily\bfseries\footnotesize, inner sep=1pt] at (-0.5,4) {2\,cm};
\end{tikzpicture}
\end{minipage}\hfill
\begin{minipage}[c]{0.58\linewidth}
この \R{形}{かたち}の \R{面積}{めん|せき}は\\\R{何}{なん}\,cm\textsuperscript{2}？ 2つに \R{分}{わ}けよう。
\end{minipage}
\bkshikiran
\bkjans{\mbox{2 × 2 ＝ 4}、 \mbox{2 × 5 ＝ 10}、 \mbox{4 ＋ 10 ＝ 14}\quad 14\,cm\textsuperscript{2}}
\end{jibun}

\clearpage
\begin{yattemiyou}[\R{新聞紙}{しん|ぶん|し}で 1\,m\textsuperscript{2} を つくろう]
\begin{enumerate}[label={\bkcircled[bkGreen]{\arabic*}}, leftmargin=2.2\zw, itemsep=2pt, topsep=0pt]
\item \R{新聞紙}{しん|ぶん|し}を テープで つないで、 1\R{辺}{へん}が 1\,m の \R{正方形}{せい|ほう|けい}を つくろう。 \R{長}{なが}さは 1\,m の ものさし（1-8）で はかるよ。
\item かどが \R{直角}{ちょっ|かく}に なって いるか、 \R{三角}{さん|かく}じょうぎで たしかめよう（1-19）。
\item できた 1\,m\textsuperscript{2} の \R{上}{うえ}に、 \R{何人}{なん|にん} \R{立}{た}てるか ためそう。
\item \R{教室}{きょう|しつ}の ゆかに しいて、 だいたい \R{何}{なん}まい\R{分}{ぶん} しけそうか、 かんがえよう。
\end{enumerate}
\begin{bkfigure}
\begin{tikzpicture}[x=1mm, y=1mm]
  \fill[white, rounded corners=3mm] (0,0) rectangle (122,50);
  \filldraw[fill=bkGray!12, draw=bkInk!60, line width=0.8pt] (14,6) -- (74,6) -- (90,30) -- (30,30) -- cycle;
  \foreach \t in {0.5} {\draw[bkGray!50, line width=0.5pt] ({14+60*\t},6) -- ({30+60*\t},30);}
  \draw[bkGray!50, line width=0.5pt] (22,18) -- (82,18);
  \bkkid{(38,12)}{24}{bkRoute6}{bkHair}\bkkid{(52,13)}{23}{bkRoute4}{bkHair}\bkkid{(66,12)}{24}{bkRoute2}{bkHair}
  \node[anchor=north, font=\sffamily\bfseries\footnotesize] at (44,5) {1\,m};
  \node[anchor=west, font=\sffamily\bfseries\footnotesize] at (83,17) {1\,m};
  \node[anchor=west, font=\sffamily\footnotesize, text=bkGray, align=left] at (94,32) {\R{新聞紙}{しん|ぶん|し}を\\つないだ\\1\,m\textsuperscript{2}};
\end{tikzpicture}
\end{bkfigure}
\noindent\hspace*{1\zw}\R{立}{た}てた \R{人数}{にん|ずう}\enspace\bkkakikomi\,\R{人}{にん}\hspace{2\zw}\R{教室}{きょう|しつ}に しける まい\R{数}{すう}\enspace だいたい\enspace\bkkakikomi[16mm]\,まい\par
\noindent\hspace*{1\zw}{\small\color{bkOrange!58!black}ヒント：\R{教室}{きょう|しつ}が たて 7\,m、 よこ 9\,m なら、 1\,m\textsuperscript{2} が 7 × 9 ＝ 63まい\R{分}{ぶん}（\R{数学}{すう|がく}の \R{道具}{どう|ぐ}④）。}
\end{yattemiyou}

\par\vspace{6pt}
%% ---------------- p.218 まちがえやすいところ ----------------
\bksection{tsumazuki}{まちがえやすい ところ}

\begin{tsumazuki}{まわりが \R{長}{なが}い ほうが \R{広}{ひろ}い？}
\begin{minipage}[c]{0.56\linewidth}
まわりの \R{長}{なが}さを マスの \R{辺}{へん}で かぞえると、 A は 14こ\R{分}{ぶん}、 B は 12こ\R{分}{ぶん}。 A の ほうが \R{長}{なが}いです。
でも \R{広}{ひろ}さは、 A が 6マス、 B が 9マス。 \R{広}{ひろ}いのは B です。
\end{minipage}\hfill
\begin{minipage}[c]{0.42\linewidth}\centering
\begin{tikzpicture}[x=1mm, y=1mm]
  \bkqtcells{(0,4)}{4.5}{bkYellow!60}{0/0,1/0,2/0,3/0,4/0,5/0}
  \bkqttag{(13.5,12)}{A}
  \bkqtcells{(36,0)}{4.5}{bkRoute4!40}{0/0,1/0,2/0,0/1,1/1,2/1,0/2,1/2,2/2}
  \bkqttag{(42.75,16.5)}{B}
\end{tikzpicture}
\end{minipage}
\end{tsumazuki}

\begin{tsumazuki}{たし\R{算}{ざん}に しない}
たて 3\,cm、 よこ 4\,cm の \R{長方形}{ちょう|ほう|けい}の \R{面積}{めん|せき}\\
\bkbatsu\enspace 3 ＋ 4 ＝ 7\hspace{3\zw}\bkmaru\enspace 3 × 4 ＝ 12
\end{tsumazuki}

\begin{tsumazuki}{1\R{辺}{へん} × 4 は まわりの \R{長}{なが}さ}
1\R{辺}{へん} 3\,cm の \R{正方形}{せい|ほう|けい}の \R{面積}{めん|せき}\\
\bkbatsu\enspace 3 × 4 ＝ 12\hspace{3\zw}\bkmaru\enspace 3 × 3 ＝ 9\par
3 × 4 は、 まわりの \R{長}{なが}さ（3\,cm の \R{辺}{へん}が 4\R{本}{ほん}）です。 \R{面積}{めん|せき}は たて × よこ。
\end{tsumazuki}

\begin{column}{はっぱの \R{広}{ひろ}さを はかろう}
\begin{minipage}[c]{0.64\linewidth}
まがった \R{形}{かたち}でも、 マスを かぞえて たし\R{合}{あ}わせれば、 だいたいの \R{広}{ひろ}さが わかります。 マスを \R{小}{ちい}さく するほど、 ほんとうの \R{広}{ひろ}さに \R{近}{ちか}く なります。
\end{minipage}\hfill
\begin{minipage}[c]{0.33\linewidth}\centering
\begin{tikzpicture}[x=1mm, y=1mm]
  \fill[bkLeaf!55] (0,10) .. controls (8,22) and (26,20) .. (36,10) .. controls (26,0) and (8,-2) .. (0,10) -- cycle;
  \draw[bkLeaf!80!black, line width=0.8pt] (-3,10) -- (36,10);
  \draw[bkBlue!45, line width=0.4pt] (-3,0) grid[step=4] (37,20);
\end{tikzpicture}
\end{minipage}
\end{column}

\begin{tsumazuki}{たんいは cm では なく cm\textsuperscript{2}}
\bkbatsu\enspace 12\,cm\hspace{3\zw}\bkmaru\enspace 12\,cm\textsuperscript{2}\par
\R{長}{なが}さは cm、 \R{広}{ひろ}さは cm\textsuperscript{2} です。
\end{tsumazuki}

\begin{column}{ひいて もとめる やり\R{方}{かた}}
\begin{minipage}[c]{0.64\linewidth}
れいだい6の \R{形}{かたち}は、 たて 4\,cm、 よこ 6\,cm の \R{長方形}{ちょう|ほう|けい}から、 \R{右上}{みぎ|うえ}の たて 2\,cm、 よこ 3\,cm を とった \R{形}{かたち}とも \R{見}{み}られます。
4 × 6 ＝ 24、 2 × 3 ＝ 6、 24 － 6 ＝ 18。 やっぱり 18\,cm\textsuperscript{2} です。
\end{minipage}\hfill
\begin{minipage}[c]{0.33\linewidth}\centering
\begin{tikzpicture}[x=1mm, y=1mm]
  \fill[bkRoute4!35] (0,0) -- (24,0) -- (24,8) -- (12,8) -- (12,16) -- (0,16) -- cycle;
  \draw[bkRed, line width=0.8pt, dash pattern=on 2pt off 1.5pt, fill=bkRed!8] (12,8) rectangle (24,16);
  \bkgrid{(0,0)}{6}{4}{4}
  \draw[bkRoute4!70!black, line width=1.2pt, line join=round] (0,0) -- (24,0) -- (24,8) -- (12,8) -- (12,16) -- (0,16) -- cycle;
  \node[font=\sffamily\bfseries\footnotesize, text=bkRed] at (18,12) {－6};
\end{tikzpicture}
\end{minipage}
\end{column}





\clearpage
%% ---------------- p.219 つかわれているところ ----------------
\begin{butsuri}
\R{面積}{めん|せき}は、 くらしの いろいろな ところで つかわれて います。
\begin{bkfigure}
\begin{tikzpicture}[x=1mm, y=1mm]
  % たたみ
  \fill[white, rounded corners=3mm] (0,0) rectangle (39,42);
  \begin{scope}[shift={(10.5,10)}]
  \bkqttatami{(0,0)}{12}{6}\bkqttatami{(12,0)}{6}{12}\bkqttatami{(0,6)}{6}{12}\bkqttatami{(6,6)}{6}{12}
  \bkqttatami{(12,12)}{6}{12}\bkqttatami{(0,18)}{12}{6}
  \draw[bkStem, line width=1.2pt] (0,0) rectangle (18,24);
  \end{scope}
  \node[font=\sffamily\bfseries\small] at (19.5,4) {6じょうの へや};
  % ソーラーパネル
  \begin{scope}[shift={(42,0)}]
  \fill[white, rounded corners=3mm] (0,0) rectangle (39,42);
  \fill[bkStem!25] (6,10) rectangle (33,20);
  \fill[bkRed!60] (3,20) -- (19.5,33) -- (36,20) -- cycle;
  \foreach \i in {0,1,2} {\foreach \j in {0,1} {\fill[bkBlue!70!black] ({10.5+5.2*\i+2.6*\j},{22+4.2*\j}) rectangle ++(4.6,3.6);}}
  \fill[bkStem!70] (16,10) rectangle (22,16);
  \node[font=\sffamily\bfseries\small] at (19.5,4) {ソーラーパネル};
  \end{scope}
  % タブレット
  \begin{scope}[shift={(84,0)}]
  \fill[white, rounded corners=3mm] (0,0) rectangle (39,42);
  \fill[bkInk!80, rounded corners=2mm] (4,11) rectangle (24,37);
  \fill[bkBlue!25] (5.5,14) rectangle (22.5,35);
  \draw[bkGray, line width=0.6pt] (17,26) circle (3.2);
  \draw[bkGray, line width=0.6pt] (19.3,23.7) -- (25.5,17.5);
  \begin{scope}
  \clip (31,15) circle (6.5);
  \foreach \i in {0,...,6} {\foreach \j in {0,...,6} {\fill[bkBlue!\the\numexpr 20+8*\i+5*\j\relax] ({24.6+1.9*\i},{8.6+1.9*\j}) rectangle ++(1.6,1.6);}}
  \end{scope}
  \draw[bkGray, line width=0.8pt] (31,15) circle (6.5);
  \node[font=\sffamily\bfseries\small] at (19.5,4) {タブレットの \R{画面}{が|めん}};
  \end{scope}
\end{tikzpicture}
\end{bkfigure}
\begin{itemize}[label={\color{bkPurple}●}, leftmargin=1.4\zw, itemsep=0pt, topsep=0pt]
\item へやの \R{広}{ひろ}さは、 たたみの \R{数}{かず}で いいます。 「6じょう」は たたみ 6まい\R{分}{ぶん}の \R{広}{ひろ}さです。
\item \R{同}{おな}じように \R{日}{ひ}が あたれば、 ソーラーパネルの \R{面積}{めん|せき}が 2つ\R{分}{ぶん}に なると、 つくれる \R{電気}{でん|き}も 2つ\R{分}{ぶん}。
\item タブレットの \R{画面}{が|めん}は、 \R{小}{ちい}さな \R{四角}{し|かく}い \R{点}{てん}の マスが ならんで できて います。
\end{itemize}
\end{butsuri}

\bkfukidashi{delta}{\R{小}{ちい}さな マスを たし\R{合}{あ}わせて \R{広}{ひろ}さを だす \R{考}{かんが}えは、\\\R{第}{だい}5\R{巻}{かん}の 「せきぶん」に つながるよ！}

\begin{bkfigure}
\begin{tikzpicture}[x=1mm, y=1mm]
  \draw[bkRoute2, line width=3pt, line cap=round] (10,0) -- (122,0);
  \node[anchor=west, inner sep=0pt] at (0,12) {\bkrouteicon[1.4em]{2}\hspace{0.3em}{\sffamily\bfseries\small\color{bkRoute2!80!black}たしあわせの \R{路線}{ろ|せん}}};
  \foreach \x/\c/\t/\s in {10/bkVol1/{1-12}/{かけ\R{算}{ざん}},38/bkVol1/{1-21}/{ひろさ},66/bkVol2/{2-3}/{{\ltjsetparameter{jacharrange={+200}}＋－×÷}},94/bkVol4/{4-5}/{こまかく たす},122/bkVol5/{5-10}/{せきぶん}} {
    \fill[white, draw=\c, line width=1.6pt] (\x,0) circle (2.6);
    \node[anchor=base, font=\sffamily\bfseries\footnotesize, text=\c!85!black] at (\x,-6.8) {\t};
    \node[anchor=base, font=\sffamily\footnotesize, text=bkInk] at (\x,-13.6) {\s};}
  \node[fill=bkTheme, text=white, rounded corners=1mm, font=\sffamily\bfseries\fontsize{8}{9}\selectfont, inner sep=1.5pt] at (38,5.5) {いま ここ};
\end{tikzpicture}
\end{bkfigure}

\clearpage
%% ---------------- p.220〜221 もんだい ----------------
\bkendproblems

\bkproblemgroup{きほん}{1}

\bkproblem A と B、 どちらが \R{何}{なん}マス \R{広}{ひろ}いですか。
\begin{bkfigure}
\begin{tikzpicture}[x=1mm, y=1mm]
  \bkgrid{(0,0)}{14}{3}{5}
  \bkqtcells{(0,0)}{5}{bkYellow!60}{0/0,1/0,2/0,3/0,0/1,1/1,2/1,3/1,0/2,1/2,2/2,3/2}
  \bkqtcells{(0,0)}{5}{bkRoute4!40}{6/0,7/0,8/0,9/0,10/0,6/1,7/1,8/1,9/1,6/2,7/2,8/2,9/2}
  \bkqttag{(10,17.5)}{A}\bkqttag{(42.5,17.5)}{B}
\end{tikzpicture}
\end{bkfigure}
\bkshikiran[10\zw]
\bkpans{13 － 12 ＝ 1\quad B が 1マス \R{広}{ひろ}い}

\bkproblem この \R{形}{かたち}の \R{面積}{めん|せき}は \R{何}{なん}\,cm\textsuperscript{2} ですか。 （1マスは 1\,cm\textsuperscript{2}）
\begin{bkfigure}
\begin{tikzpicture}[x=1mm, y=1mm]
  \bkgrid{(0,0)}{7}{3}{5}
  \bkqtcells{(0,0)}{5}{bkRoute4!45}{1/2,2/2,3/2,4/2,5/2,2/1,3/1,4/1,2/0,3/0,4/0}
\end{tikzpicture}
\end{bkfigure}
\bkkotaeran
\bkpans{11\,cm\textsuperscript{2}}

\bkproblem たて 5\,cm、 よこ 7\,cm の \R{長方形}{ちょう|ほう|けい}の \R{面積}{めん|せき}は \R{何}{なん}\,cm\textsuperscript{2}？
\bkshikiran
\bkpans{5 × 7 ＝ 35\quad 35\,cm\textsuperscript{2}}

\bkproblem この \R{形}{かたち}の \R{面積}{めん|せき}を もとめましょう。
\begin{bkfigure}
\begin{tikzpicture}[x=1mm, y=1mm]
  \fill[bkRoute6!30] (0,0) -- (28,0) -- (28,8) -- (12,8) -- (12,20) -- (0,20) -- cycle;
  \bkgrid{(0,0)}{7}{5}{4}
  \draw[bkRoute6!70!black, line width=1.2pt, line join=round] (0,0) -- (28,0) -- (28,8) -- (12,8) -- (12,20) -- (0,20) -- cycle;
  \node[anchor=south, font=\sffamily\bfseries\footnotesize, inner sep=1pt] at (6,20.3) {3\,cm};
  \node[anchor=north, font=\sffamily\bfseries\footnotesize, inner sep=1pt] at (14,-0.3) {7\,cm};
  \node[anchor=east, font=\sffamily\bfseries\footnotesize, inner sep=1pt] at (-0.5,14) {3\,cm};
  \node[anchor=east, font=\sffamily\bfseries\footnotesize, inner sep=1pt] at (-0.5,4) {2\,cm};
\end{tikzpicture}
\end{bkfigure}
\bkshikiran
\bkpans{\mbox{3 × 3 ＝ 9}、 \mbox{2 × 7 ＝ 14}、 \mbox{9 ＋ 14 ＝ 23}（1-7の ひっ\R{算}{ざん}）\quad 23\,cm\textsuperscript{2}}

\bkproblemgroup{チャレンジ}{1}

\bkproblem \R{教室}{きょう|しつ}の ゆかは、 たて 7\,m、 よこ 9\,m の \R{長方形}{ちょう|ほう|けい}です。\\\R{面積}{めん|せき}は \R{何}{なん}\,m\textsuperscript{2}？
\bkshikiran
\bkpans{7 × 9 ＝ 63\quad 63\,m\textsuperscript{2}}

\bkproblem たて 4\,cm、 よこ 4\,cm の \R{正方形}{せい|ほう|けい}と、 たて 2\,cm、 よこ 8\,cm の \R{長方形}{ちょう|ほう|けい}。 どちらが \R{広}{ひろ}いですか。
\bkkotaeran[16\zw]
\bkpans{\mbox{4 × 4 ＝ 16}、 \mbox{2 × 8 ＝ 16}\quad どちらも 16\,cm\textsuperscript{2} で \R{同}{おな}じ}

\begin{furikaeri}
\begin{checklist}
\item \R{広}{ひろ}さを、 マスの \R{数}{かず}で くらべられる。\hfill{\footnotesize\sffamily\color{bkTheme}→\pageref{bk:1-21:masu}ページ}
\item 1\,cm\textsuperscript{2} を もとに、 \R{面積}{めん|せき}を あらわせる。\hfill{\footnotesize\sffamily\color{bkTheme}→\pageref{bk:1-21:cm}ページ}
\item \R{長方形}{ちょう|ほう|けい}の \R{面積}{めん|せき}を もとめられる。\hfill{\footnotesize\sffamily\color{bkTheme}→\pageref{bk:1-21:kakezan}ページ}
\end{checklist}
チェックが つかなかったら、 →の ページに もどろう。
\end{furikaeri}

\bkfukidashi{sora}{\R{第}{だい}1\R{巻}{かん}、 おつかれさま！\\かぞえる・はかる \R{道具}{どう|ぐ}が、 ぜんぶ そろったね。\\\R{第}{だい}2\R{巻}{かん}で、 また いっしょに たんけんしよう！}
\clearpage
